% Bateria final de testes da dissertação: o que é testado, como, por quê, e o que saiu até agora.
% Compilar a partir desta pasta:  latexmk -pdf bateria_de_testes.tex
% As tabelas \input{...} vêm de ../resultados/sintese (regeneradas pela bateria a cada síntese).
\documentclass[11pt,a4paper]{article}
\usepackage[utf8]{inputenc}
\usepackage[T1]{fontenc}
\usepackage[brazil]{babel}
\usepackage{lmodern}
\usepackage[margin=2.3cm]{geometry}
\usepackage{booktabs}
\usepackage{graphicx}
\usepackage{amsmath,amssymb}
\usepackage{array}
\usepackage{tabularx}
\usepackage{xcolor}
\usepackage{enumitem}
\usepackage{float}
\usepackage{lscape}
\usepackage[hidelinks]{hyperref}
\setlist{nosep,leftmargin=*}
\newcommand{\cmd}[1]{\texttt{\detokenize{#1}}}
\newcommand{\sint}{../resultados/sintese}
\newcommand{\tabelagerada}[2]{%
  \IfFileExists{#1}{\resizebox{#2}{!}{\input{#1}}}{\fbox{\parbox{0.9\linewidth}{\centering\small
  Tabela ainda não gerada: \texttt{\detokenize{#1}}}}}}
\newcommand{\figurase}[2]{%
  \IfFileExists{#1}{\includegraphics[width=#2]{#1}}{\fbox{\parbox{0.9\linewidth}{\centering\small
  Figura ainda não gerada: \texttt{\detokenize{#1}}}}}}

\title{Bateria final de testes\\[2pt]\large Fusão de odometria visual monocular e odometria de rodas em grafo de fatores,\\ com fechamento de laço}
\author{Documento de trabalho da dissertação}
\date{Estado em 9 de novembro de 2026}

\begin{document}
\maketitle

\begin{abstract}
Este documento descreve a bateria de testes que gera todos os números da dissertação: o sistema avaliado, os dados, as
métricas, cada experimento (ablação, repetições, teste cego, detector de laço isolado, robustez a laços falsos, laços de
ponta a ponta, sensibilidade a erros dos sensores, sensibilidade da porta de deriva e comparação com o ORB-SLAM3), a forma
de execução (orquestrador retomável, agendado à noite) e os resultados obtidos até a data. As tabelas de ablação, de
sensibilidade e de laços de ponta a ponta são importadas diretamente dos arquivos gerados pela bateria e se atualizam a
cada síntese; as demais trazem os números da data acima.
\end{abstract}

\tableofcontents
\newpage

%=====================================================================================================================
\section{Objetivo}

A bateria responde a quatro perguntas:
\begin{enumerate}
  \item \textbf{Quanto o sistema melhora} em relação a cada sensor sozinho (só roda, só odometria visual)?
  \item \textbf{De onde vem o ganho}: qual a contribuição de cada componente (fechamento de laço, realimentação de
        escala, âncora pelo chão, marginalização no BA, camadas de rejeição de laço)?
  \item \textbf{O resultado generaliza}? Sequências fora do ajuste (teste cego), repetições do sistema não
        determinístico, perturbações de calibração e laços falsos adversários.
  \item \textbf{Como se compara ao estado da arte} (ORB-SLAM3 monocular) no mesmo protocolo?
\end{enumerate}

Antes da bateria, os resultados vinham de execuções avulsas feitas com versões diferentes do código, o que dificultava a
comparação. A bateria fixa uma versão do código (registrada por commit no índice \cmd{resultados/README.md}), roda tudo
com ela e arquiva os resultados anteriores em \cmd{resultados/_arquivo/2026-11-07_antes_da_bateria/}.

%=====================================================================================================================
\section{Sistema avaliado}
\label{sec:sistema}

\subsection{Visão geral}
\begin{itemize}
  \item \textbf{Odometria visual (VO) monocular}: rastreamento KLT com checagem \emph{forward-backward}, keyframes,
        triangulação e BA local em janela deslizante (8 keyframes). Dois modos de \emph{gauge} do BA: fixar os 2
        keyframes mais antigos (\emph{fix}, padrão) ou marginalizar o keyframe que sai (complemento de Schur com FEJ,
        \emph{marg}). A escala monocular é ancorada pela altura conhecida da câmera sobre o chão (pontos do plano do
        chão) e, opcionalmente, realimentada pela escala estimada no grafo.
  \item \textbf{Odometria de rodas}: modelo diferencial plano a partir dos encoders (KAIST) ou pose do robô (Inomi).
        É a única fonte de escala métrica do sistema.
  \item \textbf{Grafo de fatores} (GTSAM, iSAM2), com um nó a cada 0,5\,s:
        fatores de VO (movimento relativo com escala $\exp(S_i)$ por nó, kernel de Huber), passeio aleatório da escala,
        fatores de roda e fatores de laço 4-DoF ($x, y, z$ e guinada, como no VINS-Mono~\cite{qin2018vins}).
        No fim, um lote (Levenberg--Marquardt) refina tudo: é a \emph{fusão suavizada}. A estimativa disponível no
        instante de cada frame é a \emph{fusão online} (causal).
  \item \textbf{Fechamento de laço}: detector por saco de palavras visuais (estilo DBoW2~\cite{galvez2012dbow}) com
        ilhas temporais e consistência em 3 consultas seguidas, verificação geométrica por matriz fundamental
        ($\geq 40$ inliers) e PnP sobre pontos 3D triangulados do keyframe candidato ($\geq 20$ inliers). Candidatos
        só entre keyframes separados por $\geq 30$\,s e $\geq 100$ keyframes.
\end{itemize}

\subsection{Rejeição de laços falsos (três camadas)}
\label{sec:rejeicao}
O esquema segue o Kimera-RPGO~\cite{rosinol2020kimera}: um teste contra a odometria, depois PCM, e GNC no lote final.

\paragraph{Porta de deriva.} Compara a pose relativa afirmada pelo laço com a pose relativa estimada pela odometria entre
os dois nós. O laço é descartado se a diferença passar de
\[
  \Delta p > 15\,\text{m} + 0{,}10\,d, \qquad \Delta\theta > 0{,}35\,\text{rad} + 0{,}0005\,\text{rad/m}\cdot d,
\]
onde $d$ é o percurso entre os dois nós. É uma adaptação do teste de consistência com a odometria do
Kimera-RPGO: no lugar da covariância da odometria (que subestima a própria deriva e faria o teste $\chi^2$ rejeitar
laços verdadeiros) usa-se um limite proporcional ao percurso. Os 10\,\% cobrem com folga a pior deriva final medida da
roda (8,6\,\%, urban30); os 15\,m absorvem laços curtos e o erro de escala local do VO. A sensibilidade a esses valores
é medida na Seção~\ref{sec:porta}.

\paragraph{PCM (\emph{Pairwise Consistency Maximization})~\cite{mangelson2018pcm}.} Para cada par de laços, compõe-se o
ciclo laço~A $\to$ odometria $\to$ laço~B $\to$ odometria e testa-se se ele fecha (teste $\chi^2$ com 4 graus de
liberdade a 99\,\%, limiar 13,28, com incerteza proporcional à distância). Entram no iSAM2 apenas os laços do maior
conjunto mutuamente consistente (clique máximo), com pelo menos 2 laços: um laço isolado não entra online.

\paragraph{GNC-TLS no lote final~\cite{yang2020gnc}.} Custo de mínimos quadrados truncados (TLS, limiar
$\chi^2_6$ a 99\,\% $=16{,}81$) sobre todos os candidatos, resolvido por não convexidade graduada: cada laço recebe um
peso que vai a 0 (rejeitado) ou 1 (aceito). Duas variantes:
\begin{itemize}
  \item \textbf{GNC novo} (padrão desde 2026-11-07): parte da solução só de odometria (sem laços) e só considera laços
        com pelo menos um parceiro consistente no PCM.
  \item \textbf{GNC antigo}: partia da estimativa online e considerava também laços isolados.
\end{itemize}

%=====================================================================================================================
\section{Dados}
\label{sec:dados}

\subsection{KAIST Complex Urban~\cite{jeong2019kaist}}
Carro em ambiente urbano denso (Coreia do Sul), câmera estéreo (usa-se só a esquerda, 1280$\times$560, 10\,Hz),
encoders nas rodas traseiras e GT por fusão de RTK-GPS, IMU de fibra óptica e SLAM com LiDAR. Poucos datasets públicos
têm câmera, encoder de roda e GT de qualidade ao mesmo tempo; esse é o motivo da escolha.

\begin{table}[H]\centering\small
\caption{Sequências do KAIST usadas. ``Revisitas'': amostras do GT a menos de 10\,m de um ponto já visitado depois de
$\geq 200$\,m rodados, no mesmo sentido ($<30^\circ$) ou oposto ($>150^\circ$).}
\label{tab:seqs}
\begin{tabular}{llrll}
\toprule
Sequência & Grupo & Percurso & Revisitas & Observação \\
\midrule
urban26 & avaliação & 3,7 km & mesmo sentido & a mais curta \\
urban27 & avaliação & 5,2 km & ambos & \\
urban30 & avaliação & 5,3 km & mesmo sentido (poucas) & trânsito denso; VO tem dificuldade de reinicializar \\
urban33 & avaliação & 7,4 km & \textbf{só sentido oposto} & também usada no treino do vocabulário \\
urban38 & avaliação & 11,2 km & ambos & calibração própria (diferente de 25--34) \\
urban39 & avaliação & 10,6 km & ambos & idem \\
\midrule
urban28 & teste cego & 11,2 km & ambos & mesmas ruas de urban38/39, outra sessão \\
urban31 & teste cego & 11,3 km & nenhuma & VO perdido em 42--47\,\% dos frames \\
urban32 & teste cego & 7,1 km & mesmo sentido & \\
urban37 & teste cego & 8,1 km & nenhuma & $\sim$78\,km/h; VO perdido em 57\,\% dos frames \\
\bottomrule
\end{tabular}
\end{table}

\textbf{Avaliação} são as sequências usadas para ajustar parâmetros durante o desenvolvimento. \textbf{Teste cego} são
sequências que nunca foram usadas no ajuste nem no treino do vocabulário de palavras visuais (treinado com urban25, 33 e
34). Ressalvas: urban28 é outra sessão nas mesmas ruas de urban38/39 (cega para o ajuste, não para o lugar), e urban33
está tanto na avaliação quanto no treino do vocabulário.

\subsection{Inomi (robô móvel)}
Robô diferencial em ambiente interno, câmera monocular e pose de odometria do próprio robô. Das 10 gravações, só
\cmd{dataset_16_16} e \cmd{dataset_16_16_n2} têm GT (percurso de $\sim$2,6\,m, sem revisitas). Servem para mostrar
que o sistema roda em outra plataforma, não para medir precisão. O detector de laço e os laços falsos não rodam nelas.

%=====================================================================================================================
\section{Métricas}
\label{sec:metricas}

\begin{description}
  \item[ATE] (\emph{Absolute Trajectory Error}): RMSE da posição depois de alinhar a trajetória ao GT, calculado com o
    evo~\cite{grupp2017evo}. A fusão e a roda são métricas e se alinham por SE(3); o VO monocular puro e o ORB-SLAM3 não
    têm escala e se alinham por Sim(3), o que lhes dá a escala do GT ``de graça''.
  \item[RPE] (\emph{Relative Pose Error}) no protocolo do KITTI~\cite{geiger2012kitti}: erro de translação (\%) e de
    rotação ($^\circ$/100\,m) em segmentos de 100, 200, \dots, 800\,m; reporta-se a mediana. Não é calculado na Inomi
    (percurso menor que 100\,m).
  \item[Deriva final] Erro de posição no último instante dividido pelo percurso (\%).
  \item[Laços aceitos / errados] Um laço é ``errado'' se a distância medida entre as duas poses difere da distância no
    GT em 5\,m ou mais. ``Grosseiro'' se o erro passa de 50\,m (lugar errado).
  \item[Cobertura] (só ORB-SLAM3) fração do percurso do GT coberta pelo maior mapa.
  \item[Médias] Ponderadas pelo comprimento do percurso, separadas para avaliação e teste cego.
\end{description}

\subsection{Não determinismo e repetições}
O VO puro é determinístico (bit a bit idêntico entre execuções). O sistema completo \textbf{não} é: o grafo roda em
outro processo, de forma assíncrona, e a escala estimada por ele realimenta o VO; pequenas diferenças de tempo mudam
quando cada informação chega e a trajetória diverge. Por isso o sistema completo roda \textbf{3 vezes} por sequência e
se reporta a mediana com o mínimo e o máximo. As demais configurações rodam uma vez, e uma diferença em relação ao
sistema completo só é considerada efeito real se cair fora da faixa [mínimo, máximo] das 3 repetições.

%=====================================================================================================================
\section{Experimentos}
\label{sec:experimentos}

\subsection{Ablação}
\label{sec:ablacao}
Cada linha desliga um componente do sistema completo (Tabela~\ref{tab:configs}). ``Só roda'', ``só VO'' e ``fusão
online'' não são execuções à parte: saem das trajetórias gravadas pelas execuções correspondentes.

\begin{table}[H]\centering\small
\caption{Configurações da ablação e o que cada uma altera (opções de \cmd{Codes/scripts/rodar_sistema.py}).}
\label{tab:configs}
\begin{tabularx}{\linewidth}{lX}
\toprule
Configuração & O que muda \\
\midrule
\cmd{completo} (3$\times$) & sistema inteiro (BA fixo, realimentação, âncora do chão, laço com as 3 camadas) \\
\cmd{sem_laco} & \cmd{--no-loop}: sem fechamento de laço \\
\cmd{sem_realimentacao} & \cmd{--no-feedback}: a escala do grafo não volta para o VO; fornece também o ``só VO, BA fixo'' \\
\cmd{sem_chao} & \cmd{--vo use_ground_scale=False}: sem âncora de escala pelo plano do chão \\
\cmd{ba_marg} & \cmd{--ba-mode marg}: BA local com marginalização \\
\cmd{vo_marg} & \cmd{--no-feedback --ba-mode marg}: fornece o ``só VO, BA marginalizado'' \\
\cmd{sem_porta} & porta de deriva desligada (limites $10^9$) \\
\cmd{sem_pcm} & PCM desligado ($\chi^2=\infty$, clique mínimo 1) \\
\cmd{sem_gnc} & sem GNC-TLS no lote final (LM só com os laços ativos do PCM) \\
\cmd{sem_rejeicao} & as três camadas desligadas \\
\bottomrule
\end{tabularx}
\end{table}

\subsection{Detector de laço isolado}
Mede o reconhecimento de lugar sem o grafo (\cmd{Codes/ferramentas/avaliar_lacos.py}). Keyframes a cada $\sim$2\,m;
cada um consulta o banco de lugares. Reporta revocação no topo (R@1 e R@10) para revisitas no mesmo sentido e no oposto,
precisão dos laços aceitos (lugar certo; pose certa) e cobertura das revisitas.

\subsection{Laços de ponta a ponta}
A partir dos laços que o detector entregou ao grafo em cada execução (\cmd{dados/lacos.csv}), conta em que camada cada
laço parou (porta, PCM/GNC) ou se foi aceito, e quantos dos aceitos estavam errados
(\cmd{Codes/ferramentas/lacos_ponta_a_ponta.py}). Responde ``o que as camadas fazem com laços reais''.

\subsection{Robustez a laços falsos}
\label{sec:falsos}
Repete só o back-end (o grafo) a partir das mensagens gravadas da execução \cmd{completo_r1}
(\cmd{FUSION_RECORD}), injetando laços que certamente estão errados (\cmd{Codes/ferramentas/lacos_falsos.py}):
\begin{itemize}
  \item o nó novo $b$ é o nó atual no instante da injeção; o nó antigo $a$ está $\geq 60$\,s antes e, no GT, a mais de
        50\,m de $b$;
  \item a medida afirmada imita um laço verdadeiro típico: $b$ de 0 a 8\,m à frente de $a$, $\sim$1\,m de lado, guinada
        $\sim\mathcal{N}(0, 3^\circ)$;
  \item \textbf{aleatórios}: $N$ laços independentes espalhados no percurso; \textbf{em grupos de 5}: $N/5$ grupos de 5
        laços mutuamente consistentes (uma rua inteira confundida com outra), o caso adversário do PCM;
  \item $N \in \{0, 10, 50, 100\}$, 2 sementes, as mesmas injeções em todas as configurações (completo, sem porta, sem
        PCM, sem GNC, sem nenhuma rejeição, GNC antigo).
\end{itemize}
Limitação: o replay não realimenta o VO com o grafo perturbado; mede só o efeito no grafo.

\subsection{Sensibilidade a erros dos sensores}
Sistema completo nas 6 sequências de avaliação com perturbações deliberadas da calibração (opções
\cmd{--enc-escala}, \cmd{--altura-delta}, \cmd{--atraso-roda}): diâmetro das rodas $\pm2$\,\%, altura da câmera
$\pm5$\,cm e atraso roda-câmera $\pm50$\,ms. Três execuções por perturbação (mediana [mín--máx]), comparadas com a
faixa das 3 repetições sem perturbação.

\subsection{Sensibilidade da porta de deriva}
\label{sec:porta}
Mesmo replay dos laços falsos, com a porta em cinco versões: padrão (15\,m + 10\,\%), 15\,m + 5\,\%, 15\,m + 20\,\%,
5\,m + 10\,\% e 30\,m + 10\,\%; $N \in \{0, 10, 50\}$, 2 modos, 2 sementes, nas 8 sequências com revisita. O caso
$N=0$ mostra o custo de apertar a porta (laços reais barrados); $N>0$ mostra o benefício.

\subsection{Comparação com o ORB-SLAM3 monocular}
\label{sec:orb}
ORB-SLAM3~\cite{campos2021orbslam3} compilado localmente (\cmd{~/sota}, sem viewer), com 2000 ORB por frame e a mesma
calibração de câmera do nosso VO (\cmd{Codes/ferramentas/orbslam3_kaist.py}). Quando o rastreamento se perde, ele cria
um mapa novo, desconexo e com outra escala; o protocolo usual é avaliar o \textbf{maior mapa}. Por isso reporta-se
sempre o ATE (Sim(3)) \textbf{junto com a cobertura}, e o nosso VO puro, o sistema completo e a roda são avaliados
\textbf{no mesmo trecho}. O ORB-SLAM3 também não é determinístico: 3 execuções por sequência.

%=====================================================================================================================
\section{Execução}
\label{sec:execucao}

\subsection{Orquestrador}
\cmd{Codes/ferramentas/bateria.py} monta a lista de tarefas (sequência $\times$ configuração) e as executa em ordem de
prioridade, 2 de cada vez (tarefas de replay e o ORB-SLAM3 ocupam as 2 vagas). É \textbf{retomável}: uma tarefa conta
como feita quando o arquivo final existe (\cmd{dados/metricas.csv} ou um marcador \cmd{.feito}); uma execução
interrompida é apagada e refeita. Uma trava (\cmd{fcntl}) impede duas baterias simultâneas.

\begin{table}[H]\centering\small
\caption{Fases, em ordem de prioridade.}
\begin{tabularx}{\linewidth}{rlX}
\toprule
& Fase & Tarefas \\
\midrule
1 & sistema completo & \cmd{completo_r1} (grava também o \cmd{.rec} do grafo para o replay) \\
2 & laço e realimentação & \cmd{sem_laco}, \cmd{sem_realimentacao}, GIFs \\
3 & repetições & \cmd{completo_r2}, \cmd{completo_r3} \\
4 & componentes & \cmd{sem_chao}, \cmd{ba_marg}, \cmd{vo_marg} \\
5 & camadas de rejeição & \cmd{sem_porta}, \cmd{sem_pcm}, \cmd{sem_gnc}, \cmd{sem_rejeicao} \\
6 & detector isolado & \cmd{detector} (só KAIST) \\
7 & laços falsos & \cmd{lacos_falsos} (só KAIST) \\
8 & estado da arte & \cmd{orbslam3_r1} (só KAIST) \\
9 & porta de deriva & \cmd{porta_deriva} (8 sequências com revisita) \\
10 & sensibilidade & \cmd{sens_*} (6 de avaliação) \\
11 & repetições do ORB-SLAM3 & \cmd{orbslam3_r2}, \cmd{orbslam3_r3} \\
\bottomrule
\end{tabularx}
\end{table}

\subsection{Agendamento}
Um \emph{timer} do systemd de usuário (\cmd{bateria-dissert.timer}) dispara a bateria todo dia às 21:00; ela não
inicia tarefas novas depois das 09:30 e continua na noite seguinte de onde parou. No fim de cada noite roda a síntese
(tabelas e figuras) e regenera o índice \cmd{resultados/README.md}. O computador precisa ficar ligado com a sessão aberta.

\subsection{Organização dos resultados}
\begin{itemize}
  \item \cmd{resultados/kaist/<seq>/<config>/}: cada execução (figuras \cmd{fusao_*.pdf}, \cmd{dados/} com
        trajetórias, laços, tempos e métricas, \cmd{evo/}); GIFs em \cmd{completo_r1/gifs/};
  \item \cmd{resultados/inomi/<seq>/<config>/}: idem;
  \item \cmd{resultados/sintese/}: \cmd{ablacao/}, \cmd{inomi/}, \cmd{detector_lacos/}, \cmd{lacos_falsos/},
        \cmd{lacos_ponta_a_ponta/}, \cmd{sensibilidade/}, \cmd{porta_deriva/}, \cmd{orbslam3/} e o resumo
        \cmd{RESUMO.md};
  \item \cmd{resultados/bateria/}: log e estado;
  \item \cmd{resultados/testes_unitarios/}: histórico dos testes unitários (51 testes, todos passando).
\end{itemize}

%=====================================================================================================================
\section{Resultados até a data}
\label{sec:resultados}

Estado em 11/11/2026: bateria concluída (250 tarefas). Os replays de laços
falsos e da porta de deriva foram refeitos na noite de 10/11 com o GNC corrigido. Das 30 execuções do ORB-SLAM3, 4 abortaram com erro interno
do próprio ORB-SLAM3 (\cmd{SO3::exp failed: NaN}) e contam como falha do método.

\subsection{Ablação}
A Tabela~\ref{tab:ablacao} (página em paisagem a seguir) traz o ATE de cada configuração em cada sequência.

\begin{landscape}
\begin{table}[p]\centering
\caption{ATE [m] por sequência (mediana das repetições; $\dagger$ = teste cego; médias ponderadas pelo percurso).
Gerada pela bateria.}
\label{tab:ablacao}
\tabelagerada{\sint/ablacao/ablacao_ate.tex}{\linewidth}
\end{table}
\end{landscape}

\begin{itemize}
  \item \textbf{O sistema completo reduz o ATE em 85\,\% em relação à roda} nas sequências de avaliação
        (134,8 $\to$ 19,8\,m em média ponderada) e em 91\,\% em relação ao VO puro. No teste cego, 67\,\%
        (197,5 $\to$ 65,0\,m). O RPE de translação cai de 2,4\,\% para 1,3\,\%.
  \item \textbf{O fechamento de laço é o fator dominante}: sem ele, 99,2\,m. A fusão sozinha (online) já fica em
        48,6\,m, sobretudo por corrigir a deriva de rumo da roda (urban38: 241 $\to$ 50\,m).
  \item \textbf{Realimentação, âncora do chão, BA marginalizado e camadas de rejeição não mudam o ATE de forma
        mensurável}: as diferenças ficam dentro da variação entre execuções (só urban38 varia de 12,9 a 38,7\,m).
        Exceção: sem a âncora do chão, caem os laços aceitos (152 $\to$ 136; cego 96 $\to$ 50) e sobem os errados
        (0 $\to$ 7; 1 $\to$ 3).
  \item \textbf{Teste cego}: funciona quando há revisita no mesmo sentido e o VO rastreia (urban28: 293 $\to$ 8,0\,m;
        urban32: 36,9 $\to$ 12,7\,m). Sem revisita, o ganho vem só da fusão (urban31: 244 $\to$ 101\,m) ou não existe
        (urban37, VO perdido em 57\,\% dos frames: igual à roda, sem divergir).
  \item \textbf{urban33 não tem ganho} (41\,m $\approx$ roda 38\,m) porque todas as revisitas são em sentido oposto,
        que a câmera frontal não reconhece~\cite{lowry2016survey,garg2018lost}.
\end{itemize}

\begin{figure}[H]\centering
\figurase{\sint/ablacao/repeticoes.png}{0.9\linewidth}
\caption{Variação entre as 3 execuções do sistema completo (fusão suavizada e online).}
\end{figure}

\paragraph{Por que as execuções variam.} Em urban38, o último laço acontece no km 10,2 de 11,2: o último quilômetro
roda sem correção e amplia a diferença entre execuções, que perdem o VO em pontos diferentes entre os km 3 e 6. Em
urban31, sem laço, o resultado é tão variável quanto o VO.

\subsection{Laços de ponta a ponta}
\begin{table}[H]\centering
\caption{Laços entregues pelo detector ao grafo no sistema completo (3 execuções $\times$ 10 sequências). Gerada pela bateria.}
\tabelagerada{\sint/lacos_ponta_a_ponta/tabela_completo.tex}{\linewidth}
\end{table}
Dos 789 laços entregues, 22 eram ruins ($\geq 5$\,m) e 4 grosseiros ($>50$\,m). \textbf{Nenhum grosseiro entrou}
(3 barrados pela porta, 1 pelo GNC). Entre os 754 aceitos, 5 ruins, com erro máximo de 11,5\,m (lugar certo visto de
longe). Custo: 18 laços bons barrados (2,3\,\%).

\subsection{Detector de laço isolado}
\begin{table}[H]\centering\small
\caption{Detector sem o grafo (sequências com revisita no mesmo sentido).}
\begin{tabular}{lrrrrrr}
\toprule
Seq. & Revisitas (mesmo/oposto) & R@1 mesmo & Aceitos & Precisão (lugar) & Precisão (pose) & Cobertura \\
\midrule
urban26 & 145 / 0 & 0,57 & 4 & 1,00 & 1,00 & 49\,\% \\
urban27 & 236 / 134 & 0,64 & 15 & 0,93 & 0,93 & 92\,\% \\
urban28$\dagger$ & 1114 / 595 & 0,73 & 77 & 0,91 & 0,87 & 84\,\% \\
urban30 & 121 / 0 & 0,31 & 6 & 1,00 & 0,83 & 92\,\% \\
urban32$\dagger$ & 100 / 0 & 0,83 & 8 & 0,88 & 0,88 & 100\,\% \\
urban38 & 1124 / 586 & 0,65 & 82 & 0,89 & 0,83 & 81\,\% \\
urban39 & 918 / 571 & 0,56 & 49 & 0,92 & 0,80 & 74\,\% \\
\bottomrule
\end{tabular}
\end{table}
Em revisitas de sentido oposto, R@1 $\approx 0$ em todas as sequências.

\subsection{Robustez a laços falsos}
\begin{table}[H]\centering\small
\caption{Laços falsos nas 10 sequências (120 execuções por configuração, 6360 laços falsos injetados).}
\label{tab:falsos}
\resizebox{\linewidth}{!}{%
\begin{tabular}{lrrrrr}
\toprule
Configuração & Barrados pela porta & Falsos aceitos (execuções) & Reais mantidos & ATE mediano & Pior ATE \\
\midrule
\textbf{Sistema completo} & 5959 (94\,\%) & \textbf{17 (5)} & 90\,\% & \textbf{12,9\,m} & 329\,m \\
sem porta & 0 & 51 (10) & 92\,\% & 12,9\,m & 365\,m \\
sem PCM & 5940 & 34 (12) & 83\,\% & 20,0\,m & 328\,m \\
sem GNC-TLS & 5959 & 82 (27) & 91\,\% & 22,9\,m & 563\,m \\
\textbf{sem nenhuma rejeição} & 0 & \textbf{5290 (103)} & 96\,\% & \textbf{527\,m} & diverge \\
GNC antigo & 5959 & 28 (10) & 92\,\% & 12,9\,m & 377\,m \\
\bottomrule
\end{tabular}}
\end{table}
\begin{itemize}
  \item Sem rejeição, o sistema aceita 83\,\% dos laços falsos e o ATE mediano vai a 527\,m (com casos de divergência
        do solver). Com as três camadas, 0,3\,\% entram e o ATE mediano fica no valor sem injeção.
  \item A porta barra 94\,\%; o GNC pega a maior parte do resto (sem ele, aceitos 17 $\to$ 82); o PCM protege os laços
        reais (83 $\to$ 90\,\% mantidos). Sem a porta, o GNC ainda segura o ATE mediano, mas o número de falsos aceitos
        triplica (17 $\to$ 51) e há 15 execuções com erro do solver.
  \item \textbf{GNC novo $\times$ antigo: o novo vence} (17 em 5 execuções contra 28 em 10; pior caso 329 contra 377\,m).
        A primeira versão do GNC novo empatava (28 aceitos) por um defeito corrigido em 10/11: partia do mínimos
        quadrados com todos os candidatos e parava quando os pesos estavam todos em zero; em urban39 rejeitava todos os
        laços reais (ATE 109\,m). Agora parte da odometria mais a clique do PCM e só para depois de recozer
        ($\mu\geq10$); em urban39 mantém os mesmos 641 de 702 laços reais que o GNC antigo (pior ATE 35\,m). Sem
        injeção (n = 0) o resultado é idêntico ao da versão anterior em todas as sequências, então a ablação não muda.
  \item O pior caso das duas versões é o mesmo: urban33, 100 laços em grupos (329--377\,m). Sem nenhum laço real
        (revisitas só em sentido oposto), grupos consistentes de laços falsos não têm com o que competir.
\end{itemize}
\begin{figure}[H]\centering
\figurase{\sint/lacos_falsos/lacos_falsos_ate.png}{0.95\linewidth}
\caption{ATE da fusão suavizada em função do número de laços falsos, por configuração.}
\end{figure}

\subsection{Sensibilidade a erros dos sensores}
\begin{table}[H]\centering
\caption{ATE [m] com cada perturbação: mediana [mín--máx] de 3 execuções ($^*$ = mediana fora da faixa das 3
execuções sem perturbação, para mais ou para menos). Gerada pela bateria.}
\tabelagerada{\sint/sensibilidade/sensibilidade.tex}{\linewidth}
\end{table}
Só o erro de escala do encoder tem efeito sistemático: com +2\,\%, o ATE piora em 4 das 6 sequências (mediana
+25\,\%); com $-2$\,\%, em 3 (+7\,\%). A roda é a única fonte de escala métrica. Altura $\pm5$\,cm e atraso
$\pm50$\,ms não pioram o sistema (variação mediana de $-9$ a $-6$\,\%, com tantos casos melhores quanto piores). Os
162\,m da urban38 com altura +5\,cm foram azar de uma execução: as outras duas deram 13 e 20\,m.

\subsection{Sensibilidade da porta de deriva}
\begin{table}[H]\centering\small
\caption{Porta de deriva variada no replay, com o GNC corrigido (8 sequências com revisita; 8 execuções sem injeção
e 64 com injeção de 10 e 50 laços falsos por configuração).}
\begin{tabular}{lrrrr}
\toprule
Porta & Reais mantidos (sem injeção) & Falsos barrados & Falsos aceitos & Execuções com falha \\
\midrule
\textbf{15\,m + 10\,\% (padrão)} & \textbf{95,8\,\%} & 92,1\,\% & 10 & 2 de 64 \\
15\,m + 5\,\% & 81,7\,\% & 97,2\,\% & 5 & 9 \\
15\,m + 20\,\% & 95,8\,\% & 81,9\,\% & 25 & 7 \\
5\,m + 10\,\% & 95,8\,\% & 92,3\,\% & 10 & 2 \\
30\,m + 10\,\% & 95,8\,\% & 91,8\,\% & 10 & 2 \\
\bottomrule
\end{tabular}
\end{table}
O termo proporcional de 10\,\% cobre a pior deriva medida da roda (8,6\,\%): com 5\,\% a porta barra laços verdadeiros
(urban30 perde todos e vai a 79\,m); com 20\,\% entram 2,5 vezes mais laços falsos e as falhas triplicam (2 $\to$ 7). Com o GNC corrigido, as falhas da
configuração padrão caíram de 7 para 2 em 64 execuções. O termo fixo (5, 15 ou
30\,m) quase não importa, porque nos laços com centenas de metros entre as passagens quem manda é o termo proporcional.

\subsection{Comparação com o ORB-SLAM3}
\begin{table}[H]\centering\small
\caption{ATE [m] no trecho coberto pelo maior mapa do ORB-SLAM3: mediana [mín--máx] das execuções que não abortaram.
Nossos métodos avaliados no mesmo trecho. $\dagger$ = teste cego.}
\resizebox{\linewidth}{!}{%
\begin{tabular}{lrrrrr}
\toprule
Seq. & Cobertura & ORB-SLAM3 (Sim(3)) & Nosso VO puro (Sim(3)) & \textbf{Nosso sistema} (SE(3)) & Abortos \\
\midrule
urban26 & 32\,\% & 10,4 [8,9--20,9] & 12,8 & \textbf{2,3} & 0/3 \\
urban27 & 100\,\% & 19,9 [17,1--41,7] & 93,2 & \textbf{12,5} & 0/3 \\
urban28$\dagger$ & 92\,\% & 207,8 [134--420] & 423,0 & \textbf{7,2} & 0/3 \\
urban30 & 48\,\% & 105,2 [47--163] & 44,1 & \textbf{7,5} & 1/3 \\
urban31$\dagger$ & 8\,\% & 88,5 [1,8--175] & 54,5 & \textbf{2,4} & 1/3 \\
urban32$\dagger$ & 26\,\% & 14,7 [8,8--41,4] & 4,6 & \textbf{2,7} & 0/3 \\
urban33 & 22\,\% & 60,1 [8,1--112] & 6,7 & \textbf{5,2} & 1/3 \\
urban37$\dagger$ & 12\,\% & 64,7 [14,7--70,3] & perdido & \textbf{1,2} & 0/3 \\
urban38 & 73\,\% & 72,3 [48--116] & 373,0 & \textbf{6,9} & 0/3 \\
urban39 & 34\,\% & 33,2 [23,3--43,2] & \textbf{9,2} & 16,5 & 1/3 \\
\bottomrule
\end{tabular}}
\end{table}
\begin{itemize}
  \item O ORB-SLAM3 monocular raramente cobre o percurso inteiro: só na urban27 o maior mapa tem 100\,\%; nas demais
        ele se fragmenta em vários mapas (8 a 92\,\%), e 4 das 30 execuções abortaram.
  \item No trecho que ele cobre, o nosso sistema tem ATE menor em 9 das 10 sequências, por escala métrica sem o GT
        (o Sim(3) dá ao ORB-SLAM3 a escala do GT), continuidade quando o VO cai e fechamento de laço funcionando.
  \item Quando não se perde, o VO do ORB-SLAM3 é melhor que o nosso VO puro (urban27: 20 contra 93\,m; urban38: 72
        contra 373\,m). A dissertação não deve afirmar que o nosso VO supera o estado da arte.
  \item A exceção é a urban39, onde o nosso VO puro (9,2\,m) supera o sistema completo (16,5\,m) no trecho de 34\,\%.
\end{itemize}
\begin{figure}[H]\centering
\figurase{../resultados/kaist/urban27/orbslam3_r1/comparacao_orbslam3.png}{\linewidth}
\caption{urban27: trajetórias no trecho do maior mapa do ORB-SLAM3 e erro ao longo do trecho.}
\end{figure}

\subsection{Inomi}
Num percurso de 2,6\,m sem laço, todas as configurações ficam entre 2 e 4\,cm (sistema completo: 3,7 e 1,8\,cm;
roda: 3,7 e 1,9\,cm). A fusão segue a roda, em que o grafo confia mais.

%=====================================================================================================================
\section{Limitações e pendências}
\begin{itemize}
  \item Revisitas em sentido oposto não são reconhecidas (câmera frontal): limitação do sensor; trabalhos futuros com
        descritores invariantes ao ponto de vista, câmera traseira ou LiDAR (Scan Context).
  \item O VO perde o rastreamento em alta velocidade (urban37) e trânsito denso (urban31, urban30).
  \item O sistema completo não é determinístico; todo resultado do sistema completo deve ser lido como mediana com faixa.
  \item O replay dos laços falsos não realimenta o VO.
  \item Pendentes fora da bateria: investigar o modo de falha do GNC novo; repetir urban38 com altura +5\,cm; novos
        conjuntos com o robô da plataforma SLAM-Teaching (importador e calibrador prontos).
\end{itemize}

%=====================================================================================================================
\section{Reprodução}
\begin{verbatim}
# estado e log
python3 Codes/ferramentas/bateria.py --estado
tail -f resultados/bateria/bateria.log
# rodar o que falta agora (sem limite de horário) / só a síntese
python3 Codes/ferramentas/bateria.py --sem-limite
python3 Codes/ferramentas/bateria.py --sintese
# uma tarefa isolada
python3 Codes/ferramentas/bateria.py --seqs urban32 --configs detector --sem-limite
# testes unitários
python3 Codes/tests/rodar_testes.py
# este documento
cd documentos && latexmk -pdf bateria_de_testes.tex
\end{verbatim}

%=====================================================================================================================
\begin{thebibliography}{99}\small
\bibitem{campos2021orbslam3} C. Campos, R. Elvira, J. J. G. Rodríguez, J. M. M. Montiel, J. D. Tardós.
  ORB-SLAM3: An Accurate Open-Source Library for Visual, Visual--Inertial and Multi-Map SLAM. \emph{IEEE T-RO}, 2021.
\bibitem{galvez2012dbow} D. Gálvez-López, J. D. Tardós. Bags of Binary Words for Fast Place Recognition in Image
  Sequences. \emph{IEEE T-RO}, 2012.
\bibitem{garg2018lost} S. Garg, N. Suenderhauf, M. Milford. LoST? Appearance-Invariant Place Recognition for Opposite
  Viewpoints using Visual Semantics. \emph{RSS}, 2018.
\bibitem{geiger2012kitti} A. Geiger, P. Lenz, R. Urtasun. Are we ready for Autonomous Driving? The KITTI Vision
  Benchmark Suite. \emph{CVPR}, 2012.
\bibitem{grupp2017evo} M. Grupp. evo: Python package for the evaluation of odometry and SLAM.
  \url{https://github.com/MichaelGrupp/evo}, 2017.
\bibitem{jeong2019kaist} J. Jeong, Y. Cho, Y.-S. Shin, H. Roh, A. Kim. Complex Urban Dataset with Multi-level Sensors
  from Highly Diverse Urban Environments. \emph{IJRR}, 2019.
\bibitem{lowry2016survey} S. Lowry et al. Visual Place Recognition: A Survey. \emph{IEEE T-RO}, 2016.
\bibitem{mangelson2018pcm} J. G. Mangelson, D. Dominic, R. M. Eustice, R. Vasudevan. Pairwise Consistent Measurement
  Set Maximization for Robust Multi-Robot Map Merging. \emph{ICRA}, 2018.
\bibitem{qin2018vins} T. Qin, P. Li, S. Shen. VINS-Mono: A Robust and Versatile Monocular Visual-Inertial State
  Estimator. \emph{IEEE T-RO}, 2018.
\bibitem{rosinol2020kimera} A. Rosinol, M. Abate, Y. Chang, L. Carlone. Kimera: an Open-Source Library for Real-Time
  Metric-Semantic Localization and Mapping. \emph{ICRA}, 2020.
\bibitem{yang2020gnc} H. Yang, P. Antonante, V. Tzoumas, L. Carlone. Graduated Non-Convexity for Robust Spatial
  Perception: From Non-Minimal Solvers to Global Outlier Rejection. \emph{IEEE RA-L}, 2020.
\end{thebibliography}

\end{document}
